\documentclass{article}

% Packages
\usepackage{amsmath}
\usepackage{amssymb}
\usepackage{amsthm}
\usepackage{mathtools}
\usepackage{graphicx}
\usepackage{hyperref}

% Document information
\title{Tropical Geometry with Applications to Combinatorial Auctions}
\author{Eduardo and Ruben}
\date{\today}

\newcommand{\NN}{\mathbb{N}_{\geq 0}}
\newcommand{\ZZ}{\mathbb{Z}}
\newcommand{\QQ}{\mathbb{Q}}
\newcommand{\RR}{\mathbb{R}}

\newtheorem{lemma}{Lemma}
\newtheorem{definition}{Definition}
\newtheorem{example}{Example}

\begin{document}

\maketitle

\section{Introduction}
This text is a collection of thoughts.
The topic is Tropical Geometry with applications to Combinatorial Auctions.


\section{Tropical Geometry}

\subsection{Brief Introduction to the Tropical Semiring}

\begin{definition}
	The tropical numbers are the set $\NN \cup \{\epsilon\}$ with the following operations:
	\begin{align*}
		a \oplus b  & = \max\{a,b\} \\
		a \otimes b & = a + b
	\end{align*}
\end{definition}

We can also define $\oplus$ as $\min$ without changing the properties of the tropical semiring.
In this text, we will use the notation $T$ to denote the tropical numbers.

\begin{lemma}
	\begin{enumerate}
		\item $a \oplus b = b \oplus a$ for all $a,b \in T$.
		\item $a \otimes b = b \otimes a$ for all $a,b \in T$.
		\item $a \oplus (b \oplus c) = (a \oplus b) \oplus c$ for all $a,b,c \in T$.
		\item $a \otimes (b \otimes c) = (a \otimes b) \otimes c$ for all $a,b,c \in T$.
		\item $a \otimes \epsilon = \epsilon \otimes a = \epsilon$ for all $a \in T$.
		\item $a \oplus \epsilon = \epsilon \oplus a = a$ for all $a \in T$.
		\item $a \otimes 0 = 0 \otimes a = \epsilon$ for all $a \in T$.
		\item $a \oplus 0 = 0 \oplus a = a$ for all $a \in T$.
	\end{enumerate}
\end{lemma}

\subsection{Polynomials}
Since $T$ is a semiring, we can define polynomials over $T$ in the usual way.

\subsubsection{The Usual Properties}
Some of the usual properties of polynomials remain preserved in the tropical semiring.
\begin{lemma}
	Let $f, g \in T[x]$.
	\begin{enumerate}
		\item $\deg(f \otimes g) = \deg(f) + \deg(g)$.
		\item The units of $T[x]$ are the units of $T$.
		\item $T[x]$ is a (semi-) integral domain.
		\item We get to define $T[X, Y]$ as $T[X][Y]$ inductively.
	\end{enumerate}
\end{lemma}

We would like to have localization in the tropical semiring. Luckily, ideals behave as we expect them.

\begin{definition}
	A subset $I$ of a semiring $S$ is called an \emph{ideal} of $T$ if $a, b \in I$, $r \in S$ implies $a + b \in I$ and $r \otimes a \in I$.
\end{definition}
\begin{lemma}
	Let $I$ be an ideal of a semiring $S$.
	Then $S/I$ is a semiring.
\end{lemma}

\begin{lemma}
	$(x^2)$ is an ideal of $T[x]$.
\end{lemma}
\begin{proof}
	Let $f, g \in (x^2)$.
	Then $f \oplus g = x^2 \otimes f' \oplus x^2 \otimes g' = x^2 \otimes (f' \oplus g') \in (x^2)$.
	Let $r \in T[x]$.
	Then $r \otimes f = r \otimes (x^2 \otimes f') = (r \otimes x^2) \otimes f' \in (x^2)$.
\end{proof}

This tells us the localization $T[x] / (x^2)$ is well-defined, giving us equivalence classes of polynomials.

\subsubsection{The Unusual Properties}
There is some funny behaviour with polynomials in the tropical semiring that we should be aware of.
\begin{enumerate}
	\item Derivatives are not well-defined.
	\item Polynomials are not 1-1 with their functional values.
	\item The graph of a polynomial is a piecewise linear function.
	\item The roots of a polynomial are the values where the graph is not differentiable (in the usual sense).
\end{enumerate}

\begin{lemma}
	Polynomials and their functional values are not 1-1.
\end{lemma}
\begin{proof}
	Consider $f = x^2 \oplus 1x \oplus 2$ and $g = x^2 \oplus 2x \oplus 2$.
	Then $f(\epsilon) = g(\epsilon) = 2$.
	Also $f(0) = f(1) = g(0) = g(1) = 2$.
	Finally, $f(a) = g(a) = a^2$ for all $a \geq 2$.

	So $f \neq g$ but $f(a) = g(a)$ for all $a \in T$.
\end{proof}

This relationship does define an equivalence relation on $T[x]$, but we won't go deeper into this until we need to.

\subsubsection{Graphs of Polynomials}
Note that $a \otimes x^b = a + bx$. So the graph of a polynomial is always a piecewise linear function.

\subsubsection{Roots of Polynomials}
The roots of a polynomial are the values $a$ where the graph of the polynomial is not differentiable.

\subsubsection{Notes on Polynomials}
There are some more nice and funny properties of polynomials over a tropical semiring.
Until we need them, those can be found here: \url{https://scholarsarchive.byu.edu/cgi/viewcontent.cgi?article=2504&context=etd}.
For example, FTA holds under certain conditions, roots can be defined, and we have Bezout's theorem.

\section{Combinatorial Auctions}
A combinatorial auction is an auction where bidders can make contingent bids on bundles of items.
A bid consists of a price and a bundle of items.
The auctioneer tries to resolve the auction by maximizing the total revenue.

\subsection{Encoding a Combinatorial Auction as a Tropical Polynomial}
We will explore multiple ways to encode a combinatorial auction as a tropical polynomial.
\begin{enumerate}
	\item A polynomial is a product of bids, where each bid is encoded as a polynomial.
	\item The auction is a polynomial, and each bid is a evaluations of the polynomial.
\end{enumerate}

The second option seems more natural, but has some problems.
For example, a bid of $3$ on items $1,2,3$ cannot necessarily be encoded as (1,1,1), since the value is fully contingent on the bundle.
The only natural way to encode this bid as an evaluation would be to encode it as either $(3,3,3)$, or $\{(3,0,0), (0,3,0), (0,0,3)\}$.
The latter removes contingency, and the former requires normalization in the polynomial.
For example, a monomial $x_1 x_2 x_3$ would have to be normalized by dividing by 3.
But we have no natural division in the tropical semiring.
So already, we have to jump through hoops to encode the bids as evaluations of a polynomial.
For this reason, we will first explore the first option.

\subsection{Encoding bids as polynomials}
In this section we formalize the combinatorial auction as a tropical polynomial.
More concretely, a tropical auction on $i$ items is a polynomial $P(x_i)$ in the localization $T[x_i] / (x_i^2)$.
The localization is necessary to prevent double allocations of items.
We can encode a combinatorial auction as follows.

\begin{definition}
	A tropical bid on a set of items $I$ with value $b$ is denoted as a polynomial $ p(x_i) = 0 \oplus b\bigotimes_{i \in I} x_i$.
\end{definition}

This has some nice properties. For example, the $\oplus$ of two bids on the same items leaves only the maximum bid.
Multiplying two bids gives the sum of the two bids, plus the product of the bids, which represents the realization of both bids.
We can easily identify invalid combinations of bids by looking at the order of the factors in the monomials.
Using these properties, we can encode a combinatorial auction as follows.

\begin{definition}
	A tropical auction is a product of tropical bids.
\end{definition}

We will now explore this encoding with some examples.

\subsubsection{Encoding the Auction as a Polynomial}

\begin{example}
	Suppose an auction with three bids:
	\begin{enumerate}
		\item Bidder 1 bids on item 1 with value 1.
		\item Bidder 2 bids on item 2 with value 1.
		\item Bidder 3 bids on item 3 with value 1.
	\end{enumerate}
	Each bid is a monomial $p_j(x_i) = 0 + 1 x_j$.
	The auction becomes the product
	\begin{align*}
		P(x_i) & = p_1(x_i) \otimes p_2(x_i) \otimes p_3(x_i)                                                                       \\
		       & = (0 \oplus 1 x_1) \otimes (0 \oplus 1 x_2) \otimes (0 \oplus 1 x_3)                                               \\
		       & = 0 \oplus 1 x_1 \oplus 1 x_2 \oplus 1 x_3 \oplus 2 x_1 x_2 \oplus 2 x_1 x_3 \oplus 2 x_2 x_3 \oplus 3 x_1 x_2 x_3 \\
	\end{align*}

	It is tempting to further simplify this polynomial to $P(x_i) = 0 \oplus 3 x_1 x_2 x_3$, since for any $(x_i) \in \NN$ we have $P(x_i) = 3 x_1 x_2 x_3$.
	It is important to remember that $\epsilon \in T$, and therefore, for example, $P(\epsilon, 1, 1) = 2$.
	So we cannot simplify the polynomial further.
	However, this polynomial does encode the full outcome space of the auction.
	We see that the maximum revenue is $P(1,1,1) = 3$.
\end{example}

\begin{example}
	Suppose an auction with three bids
	\begin{enumerate}
		\item Bidder 1 bids on item 1 with value 1.
		\item Bidder 2 bids on item 1, 2 with value 2.
		\item Bidder 3 bids on item 2, 3 with value 2.
	\end{enumerate}
	The bids are
	\begin{align*}
		p_1(x_i) & = 0 \oplus 1 x_1              \\
		p_2(x_i) & = 0 \oplus 2 x_1 \oplus 2 x_2 \\
		p_3(x_i) & = 0 \oplus 2 x_2 \oplus 2 x_3
	\end{align*}

	The auction becomes the product
	\begin{align*}
		P(x_i) & = p_1(x_i) \otimes p_2(x_i) \otimes p_3(x_i)                                                                      \\
		       & = (0 \oplus 1 x_1) \otimes (0 \oplus 2 x_1 \oplus 2 x_2) \otimes (0 \oplus 2 x_2 \oplus 2 x_3)                    \\
		       & = 0 \oplus 1 x_1 \oplus 2 x_1 x_2 \oplus 2 x_2 x_3 \oplus 2 x_1^2 x_2 \oplus 3 x_1 x_2 x_3 \oplus 4 x_1 x_2^2 x_3
	\end{align*}

	Each higher order term represents a double allocation of an item.
	These higher order terms vanish under localization.
	\begin{align*}
		P(x_i) & = 0 \oplus 1 x_1 \oplus 2 x_1 x_2 \oplus 2 x_2 x_3 \oplus 3 x_1 x_2 x_3
	\end{align*}

	We see that the maximum revenue is $P(1,1,1) = 3$.
\end{example}

\subsubsection{Updating the Auction Polynomial}
This example shows that the polynomial encoding of the auction can be updated to reflect new bids.

\begin{example}
	Suppose we have the same auction as in the previous example, with polynomial $P(x_i) = 0 \oplus 1 x_1 \oplus 2 x_1 x_2 \oplus 2 x_2 x_3 \oplus 3 x_1 x_2 x_3$.
	Now a new bid comes in from bidder 4, who bids on item 1 with value 2.
	We compute the new polynomial
	\begin{align*}
		P'(x_i) & = P(x_i) \otimes (0 \oplus 2 x_1)                                                                                                                                  \\
		        & = (0 \oplus 1 x_1 \oplus 2 x_1 x_2 \oplus 2 x_2 x_3 \oplus 3 x_1 x_2 x_3) \otimes (0 \oplus 2 x_1)                                                                 \\
		        & = 0 \oplus 1 x_1 \oplus 2 x_1 x_2 \oplus 2 x_2 x_3 \oplus 3 x_1 x_2 x_3 \oplus 2 x_1 \oplus 3 x_1^2 \oplus 3 x_1^2 x_2 \oplus 4 x_1 x_2 x_3 \oplus 5 x_1^2 x_2 x_3 \\
		        & = 0 \oplus 2 x_1 \oplus 2 x_1 x_2 \oplus 2 x_2 x_3 \oplus 4 x_1 x_2 x_3                                                                                            \\
	\end{align*}
\end{example}

\subsubsection{Removing Items from the Auction}
This example shows that the polynomial encoding supports the removal of items from the auction.

\begin{example}
	Suppose we have the same updated auction as in the previous example.
	Suppose for some reason, we decide that item 1 is no longer available.
	We can compute the new auction polynomial by setting $x_1 = \epsilon$.
	\begin{align*}
		P'(x_i | x_1 = \epsilon) & = 0 \oplus 2 \epsilon \oplus 2 \epsilon x_2 \oplus 2 x_2 x_3 \oplus 3 \epsilon x_2 x_3 \oplus 4 \epsilon x_2 x_3 \\
		                         & = 0 \oplus 2 x_2 x_3
	\end{align*}
\end{example}

\subsubsection{Notes on the Auction Polynomial}
There are some drawbacks to this encoding.
It is not clear how to recover the bids from the polynomial.
For example, the polynomial $P(x_i) = 0 \oplus 3 x_1 x_2 x_3$ does not tell us if the this polynomial was the result of
\begin{enumerate}
	\item One single bid on items 1, 2, 3 with value 3.
	\item Three bids on items 1, 2, 3 with value 1.
	\item One bid on item 1 with value 1 and one bid on items 2, 3 with value 2.
	\item etc.
\end{enumerate}

It remains necessary to keep track of the bids seperately.

\end{document}
